Now that we have a foundation for factorization in a ring setting, we shift our focus to polynomial rings.

\begin{recall}
Let $R$ be an integral domain. Then
\begin{enumerate}
    \item $\forall p,q \in R[x]$ with $p\neq 0$ and $q \neq 0$, $\deg(p(x)q(x)) = \deg(p(x)) + \deg(q(x))$.
    \item The units in $R[x]$ are just units of $R$.
    \item $R[x]$ is an integral domain.
\end{enumerate}
\end{recall}

If $R[x]$ is an integral domain, the field of fractions of $R[x]$ is
$$R(x) = \set{\dfrac{p(x)}{q(x)} : p,q \in R[x], q(x) \neq 0}$$

\begin{proposition}
    Let $I$ be an ideal of $R$ and let $(I) = I[x]$ (the ideal generated by $I$ in $R[x]$). Then
    $$\dfrac{R[x]}{(I)} \cong \left(\dfrac{R}{I}\right)[x]$$
    In particular, if $I$ is a prime ideal of $R$, then this is a prime ideal of $R[x]$.
\end{proposition}

\begin{proof}
    Define $\varphi : R[x] \to (R/I)[x]$ such that
    $$\sum_{i=0}^{n}a_ix^i \mapsto \sum_{i=0}^{n}\bar{a_i}x^i$$
    where $\bar{a_i}=a_i + I$. By the definition of multiplication and addition of cosets, $\varphi$ is a ring homomorphism. Let $p(x) \in R[x]$ such that 
    \begin{align*}
        p(x) &= \sum_{i=0}^{n}a_ix^i \\
        \bar{p(x)} &= \sum_{i=0}^{n}\bar{a_i}x^i
    \end{align*}
    We know 
    \begin{align*}
        p(x) \in \ker \varphi &\iff \varphi(p(x)) = \bar{0} \\
        &\iff \bar{p(x)} = \bar{0} \\
        &\iff \bar{a_i} = \bar{0}~~\forall i \\
        &\iff a_i \in I ~~\forall i
    \end{align*}
    Thus $\ker\varphi = I[x] = (I)$. Hence the result follows from the first isomorphism theorem. \\
    The second statement could be shown by letting $I$ be a prime ideal in $R$. It follows that $R/I$ is an integral domain. Hence $(R/I)[x]$ is an integral domain. Since 
    $$\left(\dfrac{R}{I}\right)[x] \cong \dfrac{R[x]}{(I)}$$
    we have that $R[x]/(I)$ is an integral domain. We conclude $(I)$ is a prime ideal.
    \qed
\end{proof}

\begin{example}
    Consider $R = \Z$ and $I = n\Z$. Then
    $$\dfrac{\Z[x]}{n\Z[x]} \cong \left(\dfrac{\Z}{n\Z}\right)[x]$$
    If $n$ is composite, then $(\Z/n\Z)[x]$ is not an integral domain. Hence, $\Z[x]/n\Z[x]$ is not an integral domain. So $n\Z[x]$ is not prime. If $n$ is a prime $p$, then 
    $$\left(\dfrac{\Z}{p\Z}\right)[x] \text{ is an integral domain.}$$
    Hence, $\Z[x]/p\Z[x]$ is an integral domain, and so $p\Z[x]$ is a prime ideal of $\Z[x]$.
\end{example}

We now focus our attention to polynomials over a field. Let $F$ be a field and define a norm on $F[x]$ by $N(p(x)) = deg(p(x)) ~~\forall p\in F[x]$.

\begin{theorem}[Polynomails Over Fields are Euclidean Domains] \leavevmode \\
    Let $F$ be a field. Then $F[x]$ is a euclidean domain. Specifically, if $a,b \in F[x]$ such that $b(x) \neq 0$ then there exists unique $q, r\in F[x]$ such that 
    $$a(x) = q(x)b(x) + r(x)$$
    where $r(x) = 0$ or $\deg (r) < \deg (b)$.
\end{theorem}

\begin{proof}
    (Existence) We proceed by induction. If $a(x) = 0$, the $q(x) = 0$ and $r(x) = 0$ gives us what we want. If $a(x)$ is a nonzero constant, then we take $q(x) = 0$ and $r(x) = a(x)$ if $\deg (b) > 0$. We have $a(x) = 0 \cdot b(x) + a(x)$. If $b(x)$ is a nonzero constant then we write $b(x) = b$ and take $q(x) = a(x)/b$, $r(x) = 0$. Then $a(x) = q(x)b(x) + r(x).$\\
    Let $\deg(a(x)) = n$ and $\deg(b(x))$ = m, and suppose the result holds for all polynomials of degree less than $n$. If $n < m$, we let $q(x) = 0$ and $a(x) = r(x)$. If $n \geq m$, let
    \begin{align*}
        a(x) &= \sum_{i=0}^{n}a_ix^i = a_nx^n + ... + a_0 \\
        b(x) &= \sum_{i=0}^{m}b_ix^i = b_mx^m + ... + b_0
    \end{align*}
    where $b_m$ and $a_n$ are nonzero. Then 
    $$a'(x) = a(x) - \dfrac{a_n}{b_m}x^{n-m} \cdot b(x)$$
    is a polynomial with $\deg(a') < \deg(a)$. By our inductive hypothesis 
    $$a'(x) = q'(x)b(x) + r(x)$$
    for some $q'(x), r(x) \in F[x]$ where $r(x) = 0$ or $\deg(r) < \deg(b)$. Letting 
    $$q(x) = q'(x) + \dfrac{a_n}{b_m}x^{n-m}$$
    we have
    \begin{align*}
        q(x)b(x)+r(x) &= \left(q'(x)+\dfrac{a_n}{b_m}x^{n-m}\right)b(x)+r(x) \\
        &= q'(x)b(x)+\dfrac{a_n}{b_m}x^{n-m}b(x) +r(x) \\
        &= a'(x) + \dfrac{a_n}{b_m}x^{n-m}b(x) \\
        &=a(x)
    \end{align*}
    completing the induction. \\
    (Uniqueness) Suppose $q_1(x)$ and $r_1(x)$ also satisfy the conditions of the theorem. Then 
    \begin{align*}
        q(x)b(x)+r(x) &= q_1(x)b(x)+r_1(x) \\
        \implies b(x)(q(x)-q_1(x)) &= r_1(x) - r(x)
    \end{align*}
    If $q(x) - q_1(x) \neq 0,$ then
    $$\deg[b(x)(q(x)-q_1(x))] = \deg(b(x)) + \deg(q(x) - q_1(x)) \geq m$$
    However, $\deg(r_1(x) - r(x)) < m$. Hence,
    \begin{align*}
        q(x) - q_1(x) &= 0 \\
        \implies q(x) &= q_1(x)
    \end{align*}
    It follows immediately that $r(x) = r_1(x)$.
    \qed
\end{proof}

\begin{corollary}
    If $F$ is a field, then $F[x]$ is a principal ideal domain and a unique factorization domain.
\end{corollary}

This result follows from the fact that all euclidean domains are principal ideal domains, and all principal ideal domains are unique factorization domains.

\begin{example}
    $\Q[x]$ is a principal ideal domain and a unique factorization domain.
\end{example}

\begin{example}
    $\Z[x]$ is not a principal ideal domain. Examine the ideal $(2,x)$ in $\Z[x]$. Recall that 
    $$(2,x) = \set{2p(x) + xq(x) : p(x),q(x) \in \Z[x]}$$
    These are polynomials with even constant coefficients.
    $$(2,x) \neq \Z[x]$$
    since $x+1 \in \Z[x]$ but not in $(2,x)$. \\
    Suppose there exists $d(x)\in \Z[x] \st (d(x)) = (2,x)$. So $2 \in (d(x))$, hence there exists $f(x) \in \Z[x]$ such that 
    $$d(x)f(x) = 2 \text{ with } \deg(2) = 0$$
    so $\deg (d(x)) = 0 = \deg (f(x))$ and 
    $$d(x) \in \set{\pm 1, \pm 2}$$
\end{example}